\chapter{Chiari Malformation}

\section*{Definition and Overview}
Chiari malformation is a structural abnormality of the base of the skull in which part of the brain -- usually the tonsils of the cerebellum -- sits lower than normal and protrudes through the foramen magnum into the upper spinal canal. This crowding can obstruct the normal flow of cerebrospinal fluid and place pressure on the brain stem and upper spinal cord. The malformation is usually present from birth, and it spans a wide spectrum: many people have no symptoms at all and are diagnosed incidentally, while others experience headaches, balance problems, weakness, and sensory disturbance. Chiari malformation is also closely associated with syringomyelia, in which a fluid-filled cavity forms within the spinal cord.

\section*{Epidemiology}
The widespread use of MRI scanning has shown that Chiari malformation is more common than was once thought, with estimates suggesting that approximately 1 in 1,000 to 1 in 5,000 people are affected. It occurs slightly more often in females than in males. Most symptomatic cases are diagnosed in adults between the ages of 25 and 45, although the condition is also found in children, sometimes in association with spina bifida. Because imaging is so often performed for unrelated reasons, many people with the malformation are diagnosed incidentally and never develop symptoms.

\section*{Aetiology and Pathophysiology}
\begin{itemize}
    \item \textbf{Congenital:} in most cases the posterior fossa (the lower part of the skull) is smaller than normal, or the hindbrain grows faster than the surrounding skull, crowding the structures at the base of the brain
    \item \textbf{Type I:} the most common form, in which the cerebellar tonsils extend into the upper spinal canal, often without significant blockage of the fluid spaces
    \item \textbf{Type II:} a more severe form nearly always associated with spina bifida, in which a greater portion of the hindbrain is displaced
    \item \textbf{Acquired (secondary):} rarely, raised pressure inside the skull from a tumour, cyst, or another cause pushes brain tissue downward through the foramen magnum
    \item \textbf{Effects:} the crowded anatomy can obstruct the flow of cerebrospinal fluid, causing it to accumulate and leading to the formation of a syrinx -- a fluid-filled cavity within the spinal cord
    \item \textbf{Genetic and developmental factors:} the exact cause is unknown in most cases, although some families show evidence of inherited susceptibility
\end{itemize}

\section*{Clinical Features}
\begin{itemize}
    \item \textbf{Headache:} typically at the back of the head or neck, often made worse by coughing, sneezing, straining, or bending forward
    \item Neck pain and stiffness
    \item Dizziness, poor balance, or clumsiness
    \item Blurred or double vision, and sensitivity to light
    \item Ringing in the ears or hearing loss
    \item Difficulty swallowing, hoarseness, or choking
    \item Numbness, tingling, or weakness in the arms or legs
    \item With a syrinx, loss of temperature sensation in the hands, and back or neck pain
    \item In some people, no symptoms at all, with the malformation found only on incidental scanning
\end{itemize}

\section*{Differential Diagnosis}
\begin{itemize}
    \item Migraine and tension-type headache
    \item Cervical disc disease and neck muscle strain
    \item Multiple sclerosis and other demyelinating conditions
    \item Syringomyelia occurring without an associated Chiari malformation
    \item Brain tumours causing raised pressure inside the skull
    \item Idiopathic intracranial hypertension
    \item Vestibular disorders causing dizziness and balance problems
    \item Cervical myelopathy from spinal cord compression
\end{itemize}

\section*{When to Seek Medical Care}
See a doctor if headaches are new, persistent, or worsened by coughing or straining, or if you develop balance problems, numbness, tingling, or weakness.

\medskip
\textbf{Call 000 (or local emergency services) for:}
\begin{itemize}
    \item Sudden severe headache unlike any previous headache
    \item New weakness or numbness in an arm or leg
    \item Difficulty swallowing or breathing
    \item Loss of coordination, seizures, or a change in consciousness
\end{itemize}

\section*{Diagnosis and Investigations}
\begin{itemize}
    \item \textbf{History:} the pattern of headaches and their triggers, associated neurological symptoms, and any history of spina bifida
    \item \textbf{Neurological examination:} testing of vision, eye movements, balance, reflexes, coordination, and sensation in the limbs
    \item \textbf{MRI of the brain and spine:} the definitive test, which shows how far the tonsils descend and whether a syrinx is present
    \item \textbf{Cine MRI:} a specialised scan that demonstrates whether cerebrospinal fluid is flowing normally around the base of the brain
    \item \textbf{CT scan:} occasionally used to assess the bones of the skull base
    \item \textbf{Sleep study:} if sleep apnoea or breathing disturbance is suspected
    \item \textbf{Referral:} to a neurologist or neurosurgeon experienced in managing Chiari malformation
\end{itemize}

\section*{Management}
\begin{itemize}
    \item \textbf{Observation:} many people with no or mild symptoms require only regular monitoring and follow-up imaging
    \item \textbf{Pain relief:} simple analgesics for headaches, and medicines to reduce intracranial pressure in selected cases
    \item \textbf{Lifestyle measures:} avoiding heavy lifting, straining, and jarring activities that provoke symptoms
    \item \textbf{Physiotherapy:} for neck pain, balance training, and muscle weakness
    \item \textbf{Surgery (posterior fossa decompression):} removal of a small amount of bone at the base of the skull to create more space for the brain and restore fluid flow
    \item \textbf{Syrinx surgery:} drainage of a large or symptomatic syrinx to prevent spinal cord damage
    \item \textbf{Multidisciplinary care:} coordinated input from neurologists, neurosurgeons, physiotherapists, and pain specialists
\end{itemize}

\section*{Complications}
\begin{itemize}
    \item Syringomyelia, with progressive weakness, numbness, and loss of sensation from the spinal cord cavity
    \item Hydrocephalus -- the build-up of fluid within the brain
    \item Sleep apnoea and other disturbances of breathing
    \item Rarely, sudden neurological deterioration from compression of the brain stem
    \item Complications of surgery, including infection, cerebrospinal fluid leak, and recurrence of symptoms
    \item Chronic headache and reduced quality of life in severe cases
\end{itemize}

\section*{Prognosis and Outcomes}
The outlook varies widely. Many people with a Chiari malformation, particularly Type I, have no symptoms and never require treatment. Among those with symptoms, surgery relieves or improves symptoms in the majority of cases, although improvement may be gradual and complete resolution is not guaranteed. People with a syrinx or the Type II form tend to have more complex courses and may need long-term specialist follow-up. Most affected people live a normal life span, and with appropriate monitoring and treatment, serious complications can usually be prevented or managed.

\section*{Prevention and Self-Care}
\begin{itemize}
    \item There is no known way to prevent the congenital form; prenatal diagnosis of spina bifida can help plan care where relevant
    \item Avoid activities that involve straining or jarring the neck if symptoms are provoked by them
    \item Practise good neck posture and take regular breaks from desk work
    \item Keep headaches and pain well managed under a doctor's guidance
    \item Attend scheduled follow-up imaging as recommended
    \item Report new weakness, numbness, or difficulty swallowing promptly
\end{itemize}

\section*{Sources and Further Reading}
The following reputable sources provide further information on Chiari malformation:
\begin{itemize}
    \item NHS. \textit{Health A--Z: conditions affecting the brain and nerves.} https://www.nhs.uk/conditions/
    \item Mayo Clinic. \textit{Diseases and conditions library.} https://www.mayoclinic.org/
    \item National Institute of Neurological Disorders and Stroke. \textit{Chiari malformation information.} https://www.ninds.nih.gov/
    \item MSD Manual. \textit{Professional edition: congenital neurological conditions.} https://www.msdmanuals.com/
    \item MedlinePlus. \textit{Health topics on brain and nervous system conditions.} https://medlineplus.gov/
\end{itemize}
